\documentclass[11pt]{article}
\usepackage[margin=25mm]{geometry}
\usepackage{amsmath,amssymb,amsthm,booktabs}
\usepackage[T1]{fontenc}
\usepackage[colorlinks=true,urlcolor=blue,citecolor=blue,linkcolor=blue]{hyperref}
\newcommand{\Tr}{\operatorname{Tr}}
\newcommand{\id}{\mathbf{1}}
\newcommand{\A}{\mathcal{A}}
\newcommand{\normone}[1]{\left\lVert #1\right\rVert_1}
\newtheorem{theorem}{Theorem}
\newtheorem{lemma}{Lemma}
\newtheorem{proposition}{Proposition}
\title{Refuting the trace-distance case of axiom (T4) of Denis, Lacaille, Martin and Serrano-Ens\'astiga}
\author{Haobo Ma\\\small ChronoAI Pte Ltd, Singapore
\and Wenlin Zhang\\\small National University of Singapore\\
\small ChronoAI Pte Ltd, Singapore}
\date{1 November 2026\\\small Licensed under
\href{https://creativecommons.org/licenses/by/4.0/}{CC BY 4.0}}
\hypersetup{pdfauthor={Haobo Ma; Wenlin Zhang},
pdftitle={Refuting the trace-distance case of axiom (T4) of Denis, Lacaille, Martin and Serrano-Ensastiga},
pdfsubject={Mixed-state anticoherence; CC BY 4.0}}
\begin{document}
\maketitle
\begin{abstract}
We refute the trace-distance case of axiom (T4) of Denis, Lacaille, Martin and
Serrano-Ens\'astiga~\cite[Sec.~III.A.2, Eq.~(22) with $p=1$; (T4) in
Sec.~II.B, Eqs.~(7)--(8)]{DLMS}, where ``a proof is left open.''
An explicit SU(2)-covariant channel decreases the order-seven anticoherence of
a mixed spin-four state by exactly $95/90552$, disproving the universal
validity suggested by their reported numerical evidence.
The channel is a normalized sum of quadrupole Kraus operations;
its covariance and trace preservation follow from Cartesian tensor identities.
Both reduced states are diagonal, so the strict inequality reduces to rational
arithmetic. We reconstruct an existing spin-four witness and give a self-contained
exact counterexample within the full mixed-state domain of the proposed axiom.
\end{abstract}

\section{The question and the witness}
Immediately after Proposition~2 in Sec.~III.A.2 of arXiv:2605.29436v1,
Denis, Lacaille, Martin and Serrano-Ens\'astiga~\cite{DLMS} report numerical
evidence for the trace-distance case of (T4) and state that
``a proof is left open.'' The proposed monotonicity is
$\A_t^{T,d_1}(\Phi(\rho))\geq\A_t^{T,d_1}(\rho)$ for every
SU(2)-covariant completely positive trace-preserving map $\Phi$
(Definition~2, Sec.~II.B, Eqs.~(7)--(8)). We refute it for a mixed state with
$j=4$ and $t=7$.
The mixed-state definition permits $1\leq t\leq N=2j$; the usual restriction
$t\leq N/2$ when discussing pure states does not restrict the mixed-state axiom.
Pure-state measures are treated in~\cite{BM}; general equivariant spin
channels are studied in~\cite{ARS}. The calculation below uses no
classification theorem.

Identify the spin-$N/2$ space with the symmetric subspace of $N$ qubits, using
the normalized Dicke basis $\{|D_N^r\rangle:0\leq r\leq N\}$,
where $r$ counts qubits in state $1$. Let $\rho_t$ denote the ordinary
partial trace over $N-t$ qubits, expressed in its normalized Dicke basis.
With the source's unhalved trace norm, define
\begin{equation}\label{eq:measure}
 \A_t^{T,d_1}(\rho)=1-
 \frac{\normone{\rho_t-\id/(t+1)}}{K_t},\qquad
 K_t=\max_{\sigma\geq0,\,\Tr\sigma=1}
 \normone{\sigma-\id/(t+1)}=\frac{2t}{t+1}.
\end{equation}
For the last equality, diagonalize $\sigma$ and express its eigenvalue vector
as a convex combination of the vertices of the probability simplex.
Convexity of the $\ell^1$ norm bounds the distance by its value at a vertex,
$2t/(t+1)$, and every vertex attains this value. Including a factor $1/2$
in both distance and $K_t$ leaves the measure unchanged.

The channel acts on the nine-dimensional irreducible spin-four space;
the eight-qubit embedding defines the reductions, not a tensor-product
noise model. Use the order $r=0,\ldots,8$, so $m=4-r$. Put
\begin{align}
 J_z|r\rangle&=(4-r)|r\rangle,&
 J_+|r\rangle&=\sqrt{r(9-r)}|r-1\rangle,\nonumber\\
 J_-&=J_+^\dagger,& J_x&=(J_++J_-)/2,\nonumber\\
 J_y&=(J_+-J_-)/(2i).&&\label{eq:spin}
\end{align}
Set $Q_{ab}=\frac12(J_aJ_b+J_bJ_a)-\frac{20}{3}\delta_{ab}\id$
for $a,b\in\{x,y,z\}$, and define
\begin{equation}\label{eq:channel}
 \Phi(X)=\frac3{770}\sum_{a,b}Q_{ab}XQ_{ab},\qquad
 \rho=\frac1{1764}\operatorname{diag}(209,92,560,0,441,0,294,168,0).
\end{equation}

\begin{theorem}\label{thm:main}
The map $\Phi$ in~\eqref{eq:channel} is completely positive, trace preserving
and SU(2)-covariant. The state $\rho$ is a mixed density operator, and
\[
 \A_7^{T,d_1}(\Phi(\rho))-\A_7^{T,d_1}(\rho)
 =-\frac{95}{90552}<0.
\]
Thus this candidate fails (T4) on the full mixed-state domain.
\end{theorem}

\section{An admissible covariant channel}
The $J_a$ are Hermitian, satisfy $[J_a,J_b]=i\sum_c\epsilon_{abc}J_c$,
and have Casimir $\sum_aJ_a^2=C\id$ with $C=20$.
Consequently every $Q_{ab}$ is Hermitian and $\sum_a Q_{aa}=0$.
The Kraus operators $\sqrt{3/770}\,Q_{ab}$ give complete positivity.
We verify their normalization rather than assuming it from covariance.

\begin{lemma}\label{lem:normalization}
For spin Casimir $C$, set $S_{ab}=(J_aJ_b+J_bJ_a)/2$ and
$Q_{ab}=S_{ab}-(C/3)\delta_{ab}\id$. Then
$\sum_{a,b}Q_{ab}^2=C(4C-3)\id/6$.
\end{lemma}
\begin{proof}
The commutator relations give
$\sum_a[J_a,[J_a,J_b]]=2J_b$. Expanding the double commutator yields
$\sum_aJ_aJ_bJ_a=(C-1)J_b$.
For $S_{ab}=(J_aJ_b+J_bJ_a)/2$, expansion of the four products gives
\[
 \sum_{a,b}S_{ab}^2
 =\tfrac12\sum_{a,b}(J_aJ_bJ_aJ_b+J_aJ_b^2J_a)
 =\tfrac12\bigl(C(C-1)+C^2\bigr)\id.
\]
Since $\sum_aS_{aa}=C\id$, subtracting the scalar trace part gives
$\sum Q_{ab}^2=(C^2-C/2-C^2/3)\id$, as asserted.
For $C=20$ this is $(770/3)\id$; hence the Kraus sum is the identity
and $\Phi$ is trace preserving. It is also unital.
\end{proof}

For any spin-four rotation $U$, there is a real orthogonal matrix $R$ with
$UJ_aU^\dagger=\sum_cR_{ac}J_c$. Therefore
\[
 UQ_{ab}U^\dagger=\sum_{c,d}R_{ac}R_{bd}Q_{cd}.
\]
Substitution into $U\Phi(X)U^\dagger$ and summation using
$\sum_aR_{ac}R_{ae}=\delta_{ce}$ gives
$U\Phi(X)U^\dagger=\Phi(UXU^\dagger)$ for every matrix $X$.
This proves covariance for every SU(2) rotation, not just a finite sample.

\section{Exact output and partial traces}
For an explicit sparse evaluation of~\eqref{eq:channel}, set $c=\sqrt{3/3080}$.
The same channel has the following five Kraus operators:
\begin{align}\label{eq:kraus}
 L_0&=(3J_z^2-20\id)/\sqrt{1540},\nonumber\\
 L_{+1}&=-c(J_zJ_++J_+J_z),&L_{-1}&=c(J_zJ_-+J_-J_z),\nonumber\\
 L_{+2}&=cJ_+^2,&L_{-2}&=cJ_-^2.
\end{align}
To check the equality of channels, use
$Q_{xx}+Q_{yy}=-Q_{zz}$ and write
$J_zJ_\pm+J_\pm J_z=2(Q_{xz}\pm iQ_{yz})$ and
$J_\pm^2=Q_{xx}-Q_{yy}\pm2iQ_{xy}$.
For Hermitian $A,B$, the cancellation identity is
\[
 (A+iB)X(A-iB)+(A-iB)X(A+iB)=2AXA+2BXB.
\]
Apply it with $(A,B)=(Q_{xz},Q_{yz})$ and
$(A,B)=(Q_{xx}-Q_{yy},2Q_{xy})$. Also, with $D=Q_{xx}-Q_{yy}$,
\[
 Q_{xx}XQ_{xx}+Q_{yy}XQ_{yy}
 =\tfrac12(DXD+Q_{zz}XQ_{zz}).
\]
The remaining
coefficients of $Q_{zz}XQ_{zz}$, $(Q_{xx}-Q_{yy})X(Q_{xx}-Q_{yy})$,
and each off-diagonal $Q_{ab}XQ_{ab}$ are respectively
$9/1540$, $3/1540$ and $3/385$, exactly those in~\eqref{eq:channel}.

For completeness, each input column of~\eqref{eq:kraus} has just one possible
output. At magnetization $m=4-r$, its squared amplitudes and output shifts are
\begin{center}\small
\begin{tabular}{ccl}
\toprule
Operator & Output index & Squared amplitude\\
\midrule
$L_0$ & $r$ & $(3m^2-20)^2/1540$\\
$L_{+1}$ & $r-1$ & $3(2m+1)^2[20-m(m+1)]/3080$\\
$L_{-1}$ & $r+1$ & $3(2m-1)^2[20-m(m-1)]/3080$\\
$L_{+2}$ & $r-2$ & $3[20-m(m+1)][20-(m+1)(m+2)]/3080$\\
$L_{-2}$ & $r+2$ & $3[20-m(m-1)][20-(m-1)(m-2)]/3080$\\
\bottomrule
\end{tabular}
\end{center}
Out-of-range amplitudes vanish. Applying these weights to the probabilities
in~\eqref{eq:channel} gives
\begin{equation}\label{eq:output}
 \Phi(\rho)=\frac1{194040}\operatorname{diag}
 (22288,30100,25528,13149,45540,10584,24654,11613,10584).
\end{equation}
The input entries are nonnegative and sum to one, with more than one nonzero
entry, proving that $\rho$ is a mixed density operator.

The normalized Dicke splitting identity is
\[
 |D_N^r\rangle=\sum_k
 \sqrt{\frac{\binom tk\binom{N-t}{r-k}}{\binom Nr}}
 |D_t^k\rangle\otimes|D_{N-t}^{r-k}\rangle.
\]
It follows by grouping computational-basis strings by their first-block
weight; the coefficient counts $\binom tk\binom{N-t}{r-k}$ strings out of
$\binom Nr$. Thus the reduction of $\operatorname{diag}(p_r)$ is diagonal with
entries
\[
 p^{(t)}_k=\sum_r p_r\frac{\binom tk\binom{N-t}{r-k}}{\binom Nr},
\]
taking binomial coefficients outside their range as zero. In particular,
for $N=8,t=7$,
\begin{equation}\label{eq:reduction}
 p^{(7)}_k=\frac{8-k}{8}p_k+\frac{k+1}{8}p_{k+1}\quad(0\leq k\leq7).
\end{equation}
Applying this formula to the input and output gives
\begin{align}\label{eq:marginals}
 \rho_7&=\frac1{168}\operatorname{diag}(21,21,40,21,21,21,21,2),\nonumber\\
 [\Phi(\rho)]_7&=\frac1{517440}\operatorname{diag}
 (69468,87252,64205,82635,78360,59892,43533,32095).
\end{align}
Since $\id/8$ corresponds to $21/168$ or $64680/517440$ in these rows,
respectively, summing absolute deviations gives
\begin{align*}
 \normone{\rho_7-\id/8}&=\frac{38}{168}=\frac{19}{84},\\
 \normone{[\Phi(\rho)]_7-\id/8}
 &=\frac{117990}{517440}=\frac{3933}{17248}.
\end{align*}
Their difference is $95/51744$. Dividing by $K_7=7/4$ in
\eqref{eq:measure} proves Theorem~\ref{thm:main}.

\section{Consequences and sibling candidates}
\paragraph{Mechanism.}
For a convex mixture of rotations, the partial trace intertwines each
collective rotation with its action on the retained qubits. Convexity and
unitary invariance of the trace norm then force the reduced distance to
decrease. The witness consequently lies outside that subclass of covariant
channels. Covariance on the original spin sector alone does not supply a
trace-norm contraction on the reduced sector.

At the full order $t=N$, by contrast, trace-distance (T4) always holds.
Indeed covariance makes $\Phi(\id/(N+1))$ invariant, hence equal to
$\id/(N+1)$ by irreducibility and trace preservation. For Hermitian $X$,
write $X=X_+-X_-$ with $X_\pm\geq0$. Positivity and trace preservation give
$\normone{\Phi(X)}\leq\Tr\Phi(X_+)+\Tr\Phi(X_-)=\normone{X}$.
Apply this with $X=\rho-\id/(N+1)$ in~\eqref{eq:measure}.

\paragraph{The convex-roof construction.}
Proposition~6 of~\cite[Sec.~III.B.2]{DLMS} assumes that the associated total
quantity satisfies the total-measure axioms. Theorem~\ref{thm:main} disproves
that hypothesis for $d_1$ on the full mixed-state domain. Consequently that
proposition cannot be applied generally to $d_1$ as a valid total--quantum
pair. This is a failure of its hypothesis, not a refutation of the conditional
proposition or a separate disproof of the quantum-side axioms.

\paragraph{Fidelity and other orders.}
The fidelity candidate in~\cite[Sec.~III.A.3, Eq.~(26)]{DLMS} is
\[
 \A_t^{T,F}(\sigma)=\frac{(\Tr\sqrt{\sigma_t})^2-1}{t}.
\]
Its (T4) remains open in that source, including the restatement in
Proposition~7 (Sec.~III.B.3). Our witness does not refute it.
\begin{proposition}
For the state and channel in~\eqref{eq:channel}, the fidelity candidate
strictly increases at every $t=1,\ldots,8$. The trace-distance candidate
strictly increases for $t=1,\ldots,6,8$ and decreases for $t=7$.
\end{proposition}
\begin{proof}
Table~\ref{tab:siblings} gives the changes, always output minus input.
\begin{table}[ht]
\centering
\begin{tabular}{rcc}
\toprule
$t$ & $\Delta\A_t^{T,d_1}$ (exact) & $10^6\Delta\A_t^{T,F}$ (lower bound)\\
\midrule
1 & $19/784$ & $3663$\\
2 & $4275/120736$ & $8574$\\
3 & $2831/75460$ & $13085$\\
4 & $4959/241472$ & $18296$\\
5 & $741/75460$ & $26268$\\
6 & $19/2464$ & $38910$\\
7 & $-95/90552$ & $62258$\\
8 & $3889/17248$ & $365375$\\
\bottomrule
\end{tabular}
\caption{Exact trace changes and certified rational fidelity bounds for this witness.}
\label{tab:siblings}
\end{table}

Here is a finite exact certificate procedure for the table. Compute every
diagonal reduction by the binomial formula above. The trace column is then
rational arithmetic. For a reduced probability $p$, set $s=10^{15}$ and
$n(p)=\lfloor\sqrt{\lfloor s^2p\rfloor}\rfloor$. Integer arithmetic gives
$n(p)^2\leq s^2p<(n(p)+1)^2$, so
$n(p)/s\leq\sqrt p<(n(p)+1)/s$.
Sum these lower and upper bounds to obtain $L_\sigma\leq
\Tr\sqrt{\sigma_t}<U_\sigma$ for each of the input and output.
Both sums are nonnegative; therefore
\[
 \Delta\A_t^{T,F}\geq
 \frac{L_{\Phi(\rho)}^2-U_\rho^2}{t}.
\]
Substitution yields a rational bound no smaller than each displayed integer
divided by $10^6$: the displayed integer is the floor of $10^6$ times this
certified bound, not a rounded estimate of the true change.
The accompanying \texttt{siblings.py} and
\texttt{siblings.json} supply the reductions and exact bounds for all eight
rows; no floating-point sign test is used.
\end{proof}

\paragraph{Remaining questions.}
The universal fidelity-based (T4) is not settled here. Nor does this calculation
settle trace-distance (T4) for mixed states restricted to $t\leq N/2$, or
trace-distance (T4) on pure inputs, in particular in the usual range
$t\leq N/2$. For arbitrary mixed-state witnesses with $1\leq t\leq2j$,
the minimum spin admitting a violation and then the minimum order at that
spin are not determined. The existence and corresponding minimum spin/order
for pure-input violations with $t\leq j$ are separate questions.
The proved purity and Hilbert--Schmidt constructions in
Propositions~1--2 of~\cite{DLMS} are not contradicted by this result.

\paragraph{Attribution.}
The state and channel are reconstructed from
\href{https://github.com/the-omega-institute/trureturing/blob/d1876cfbbef8c13023311126691b180eefa69569/D5/S3/Quantum/Measurement/AnticoherenceTraceDistanceRefutation.lean}{the pinned upstream Lean source},
where the declaration \texttt{result} states the failure of the mixed-state
claim. The argument here supplies a direct Cartesian-tensor proof and exact
reduced-state calculation of that witness, rather than a different example.

\begin{thebibliography}{3}
\bibitem{DLMS}
J. Denis, T. Lacaille, J. Martin and E. Serrano-Ens\'astiga,
\emph{Total, quantum, and classical measures of anticoherence for mixed spin
states}, arXiv:2605.29436v1 (2026), Secs.~II.B, III.A.2--3 and III.B.2--3.
\href{https://doi.org/10.48550/arXiv.2605.29436}{doi:10.48550/arXiv.2605.29436}.
\bibitem{BM}
D. Baguette and J. Martin, \emph{Anticoherence measures for pure spin states},
Phys. Rev. A \textbf{96}, 032304 (2017).
\href{https://doi.org/10.1103/PhysRevA.96.032304}{doi:10.1103/PhysRevA.96.032304}.
\bibitem{ARS}
T. Aschieri, B. Ruba and J. P. Solovej,
\emph{SU(2)-equivariant quantum channels: semiclassical analysis},
Commun. Math. Phys. \textbf{405}, 298 (2024).
\href{https://doi.org/10.1007/s00220-024-05178-1}{doi:10.1007/s00220-024-05178-1}.
\end{thebibliography}
\end{document}
